\section{Inserting Observable Kernels into the 3D Power Spectrum}

\subsection{Overview: From General Formalism to Observable Combinations}

The connection between observable-specific kernel functions and the 3D power spectrum proceeds through substitution of the window functions $\mathcal{W}_l^A(k_1, k)$ and $\mathcal{W}_l^B(k_2, k)$ in the spherical Fourier-Bessel (SFB) formula. We start with the general framework and show explicitly how each observable kernel (Eq.~\ref{obs}) enters the power spectrum computation. This derivation covers density, lensing, and their cross-correlations.

\subsection{Starting Point: The General SFB Power Spectrum}

From the main derivation, we have the general 3D power spectrum in Fourier-Bessel basis:
\begin{equation}
\label{eq:sfb_general}
S_l^{AB}(k_1,k_2) = \left(\frac{2}{\pi}\right)^2 \mathcal{N}_l^{AB} \int dk \, k^2 \, P^{AB}(k) \, \mathcal{W}_l^A(k_1, k) \, \mathcal{W}_l^B(k_2, k)
\end{equation}

where the Fourier-Bessel transform of the window functions is:
\begin{equation}
\label{eq:window_fb_transform}
\mathcal{W}_l^{A}(k, k') \equiv \int_0^{\infty} d\chi \, \widetilde{W}^{A}(k, \chi) \, j_l(k' \chi)
\end{equation}

and the redshift-space window function incorporating the spherical Bessel function is:
\begin{equation}
\label{eq:window_weighted}
\widetilde{W}^{A}(k, \chi) \equiv W^{A}_l(k, z(\chi)) \, j_l(k \chi)
\end{equation}

Here, $W^{A}_l(k, z)$ represents the observable-specific radial kernel function at wavenumber $k$ and redshift $z$. The power spectrum $P^{AB}(k)$ is the 3D matter or matter-related power spectrum at wavenumber $k$, and $\mathcal{N}_l^{AB}$ is the appropriate normalization factor for the observable pair $(A, B)$.

\subsection{Explicit Kernel Functions for Different Observables}

\subsubsection{Density Kernel (Galaxy Number Counts)}

The density kernel for galaxy clustering is:
\begin{equation}
\label{eq:kernel_den}
W^{\rm den}_l(k, z) = \frac{H(z)}{c} \, b(z) \, n(z) \, D(z) \, j_l[k r(z)]
\end{equation}

where:
\begin{itemize}
  \item $H(z) = H_0 \, E(z)$ is the Hubble parameter at redshift $z$
  \item $b(z)$ is the linear bias factor encoding how galaxy positions trace the dark matter
  \item $n(z)$ is the comoving number density of galaxies at redshift $z$
  \item $D(z)$ is the linear growth factor (normalized to $D(z=0) = 1$)
  \item $r(z)$ is the comoving distance to redshift $z$
  \item $j_l[\cdot]$ denotes the spherical Bessel function of order $l$
\end{itemize}

\textbf{Physical interpretation}: This kernel weights the dark matter fluctuations by the galaxy bias and growth, then projects via the spherical Bessel function. The $(H/c) \, n$ factor accounts for the geometry and number of tracers in each redshift shell.

\subsubsection{Weak Lensing Kernel (Convergence)}

The weak lensing convergence kernel is:
\begin{equation}
\label{eq:kernel_gamma}
W^\gamma_l(k, z) = -\sqrt{\frac{(l+2)!}{(l-2)!}} \, \frac{3}{2} \, H_0^2 \, \Omega_{{\rm m},0} \, (1+z) \, q_{3/5}(z) \, D(z) \, \frac{j_l[k r(z)]}{k^2 r(z)}
\end{equation}

where:
\begin{itemize}
  \item $\Omega_{{\rm m},0}$ is the matter density parameter at $z=0$
  \item $q_{3/5}(z)$ is the lensing efficiency function, defined as:
  \begin{equation}
  q_{3/5}(z) = \int_{r(z)}^\infty d r \, \frac{r - r(z)}{r} \, n(r) \, [5 s(r) - 2]
  \end{equation}
  accounting for the integrated lensing along the light path beyond redshift $z$
  \item The factor $\sqrt{(l+2)!/(l-2)!}$ arises from tensor projections of the lensing distortion onto the dipole component of density fluctuations
  \item The $1/(k^2 r)$ prefactor reflects the scalar nature of convergence (proportional to the second derivatives of the potential)
\end{itemize}

\textbf{Physical interpretation}: Weak lensing is sensitive to the integrated matter fluctuations along the line of sight. The kernel is weighted by the lensing efficiency and suppressed at small scales (high $k$) due to the $1/(k^2 r)$ factor.

\subsubsection{Redshift-Space Distortions Kernel}

The redshift-space distortions (RSD) kernel couples density and peculiar velocity effects:
\begin{equation}
\label{eq:kernel_rsd}
W^{\rm RSD}_l(k,z) = \frac{H(z)}{c} \, f(z) \, n(z) \, D(z) \, \left\{ \left[\frac{l(l-1)}{k^2 r^2(z)} - 1\right] j_l[k r(z)] + \frac{2}{k r(z)} j_{l+1}[k r(z)] \right\}
\end{equation}

Equivalently, in terms of the density kernel:
\begin{equation}
\label{eq:kernel_rsd_alt}
W^{\rm RSD}_l(k,z) = \frac{f(z)}{b(z)} \left\{ \left[\frac{l(l-1)}{k^2 r^2(z)} - 1\right] W^{\rm den}_l(k,z) + \frac{2}{k r(z)} W^{\rm den}_{l+1}(k,z) \right\}
\end{equation}

where:
\begin{itemize}
  \item $f(z) = d \ln D / d \ln a$ is the linear growth rate
  \item The first term represents the monopole RSD effect (density boost)
  \item The second term couples to the next Legendre multipole (quadrupole effect from velocity)
\end{itemize}

\textbf{Physical interpretation}: RSD kernels couple density and velocity fields, creating scale-dependent and scale-independent effects. The coupling between multipoles (second term) is essential for capturing the full RSD signature in 3D power spectra.

\subsubsection{Magnification Kernel}

The magnification bias kernel is related to lensing:
\begin{equation}
\label{eq:kernel_mag}
W^{\rm mag}_l(k, z) = -\sqrt{\frac{(l+2)!}{(l-2)!}} \, 3 \, H_0^2 \, \Omega_{{\rm m},0} \, (1+z) \, q_s(z) \, D(z) \, \frac{j_l[k r(z)]}{k^2 r(z)}
\end{equation}

where $q_s(z)$ is the magnification-bias weighted lensing efficiency:
\begin{equation}
q_s(z) = \int_{r(z)}^\infty d r \, \frac{r - r(z)}{r} \, n(r) \, [5 s(r) - 2]
\end{equation}

with $s(z)$ being the magnification bias parameter. The relation to the convergence kernel is:
\begin{equation}
\label{eq:mag_gamma_relation}
W^{\rm mag}_l(k,z) = 2 \, \frac{q_s(z)}{q_{3/5}(z)} \, W^\gamma_l(k,z)
\end{equation}

\subsection{Constructing the Fourier-Bessel Transform of Observable Kernels}

To implement Eq.~\ref{eq:sfb_general}, we must compute the Fourier-Bessel transform of each kernel. Consider a generic observable kernel $W_l^A(k, z)$. The steps are:

\subsubsection{Step 1: Assemble the Redshift-Integrated Window}

Start with the window function weighted by the spherical Bessel function:
\begin{equation}
\label{eq:window_construction}
\widetilde{W}^A(k, \chi) = W_l^A(k, z(\chi)) \, j_l(k \chi)
\end{equation}

This combines the observable-specific kernel (which depends on $z$ through the comoving distance $\chi = r(z)$, bias, growth, etc.) with the spatial Bessel projection.

\subsubsection{Step 2: Perform the Fourier-Bessel Transform}

The Fourier-Bessel transform over the radial direction is:
\begin{equation}
\label{eq:fb_transform_integral}
\mathcal{W}_l^A(k_1, k) = \int_0^\infty d\chi \, \widetilde{W}^A(k, \chi) \, j_l(k_1 \chi)
\end{equation}

Substituting Eq.~\ref{eq:window_construction}:
\begin{equation}
\label{eq:fb_transform_explicit}
\mathcal{W}_l^A(k_1, k) = \int_0^\infty d\chi \, W_l^A(k, z(\chi)) \, j_l(k \chi) \, j_l(k_1 \chi)
\end{equation}

This integral couples wavenumbers $k$ and $k_1$ through the redshift-dependent observable kernel. Numerically, this is typically evaluated using adaptive quadrature (e.g., Gauss-Legendre or Tanh-Sinh integration).

\subsection{Power Spectrum Components for Specific Observable Pairs}

\subsubsection{Density × Density Power Spectrum: $P^{\rm den}_l(k_1, k_2)$}

For galaxy clustering, substituting $A = B = \rm den$ into Eq.~\ref{eq:sfb_general}:
\begin{equation}
\label{eq:ps_den_den}
S_l^{\rm den,den}(k_1, k_2) = \left(\frac{2}{\pi}\right)^2 \mathcal{N}_l^{\rm den,den} \int dk \, k^2 \, P_{\rm m}(k) \, \mathcal{W}_l^{\rm den}(k_1, k) \, \mathcal{W}_l^{\rm den}(k_2, k)
\end{equation}

where $P_{\rm m}(k)$ is the matter power spectrum. The window functions are:
\begin{align}
\label{eq:window_den_fb}
\mathcal{W}_l^{\rm den}(k_1, k) &= \int_0^\infty d\chi \, \frac{H(\chi)}{c} \, b(\chi) \, n(\chi) \, D(\chi) \, j_l(k \chi) \, j_l(k_1 \chi)
\end{align}

\textbf{Normalization}: The factor $\mathcal{N}_l^{\rm den,den}$ normalizes the observable. For galaxy clustering:
\begin{equation}
\mathcal{N}_l^{\rm den,den} = 1
\end{equation}

\textbf{Physical interpretation}: This power spectrum contains the clustering of galaxies as traced through the matter density field. The double Bessel projection captures all multipole contributions beyond the Limber approximation.

\subsubsection{Weak Lensing × Weak Lensing Power Spectrum: $P^{\gamma \gamma}_l(k_1, k_2)$}

For convergence (weak lensing) auto-spectrum:
\begin{equation}
\label{eq:ps_gam_gam}
S_l^{\gamma, \gamma}(k_1, k_2) = \left(\frac{2}{\pi}\right)^2 \mathcal{N}_l^{\gamma,\gamma} \int dk \, k^2 \, P_{\rm m}(k) \, \mathcal{W}_l^\gamma(k_1, k) \, \mathcal{W}_l^\gamma(k_2, k)
\end{equation}

The lensing window functions are:
\begin{align}
\label{eq:window_gam_fb}
\mathcal{W}_l^\gamma(k_1, k) &= \int_0^\infty d\chi \, \left[-\sqrt{\frac{(l+2)!}{(l-2)!}} \, \frac{3}{2} \, H_0^2 \, \Omega_{{\rm m},0} \, (1+z(\chi)) \, q_{3/5}(\chi) \, D(\chi)\right] \\
&\quad \times \frac{j_l(k \chi)}{k^2 \chi} \, j_l(k_1 \chi)
\end{align}

\textbf{Normalization}: For convergence:
\begin{equation}
\mathcal{N}_l^{\gamma,\gamma} = \frac{1}{\pi}
\end{equation}

\textbf{Physical interpretation}: Weak lensing is more sensitive to large-scale structures due to the integrated nature of the kernel. The $1/(k^2)$ suppression at high wavenumber reflects the dominance of integrated effects over local density fluctuations.

\subsubsection{Density × Weak Lensing Cross-Spectrum: $P^{\rm den,\gamma}_l(k_1, k_2)$}

Cross-correlations between density and lensing observables yield:
\begin{equation}
\label{eq:ps_den_gam}
S_l^{\rm den,\gamma}(k_1, k_2) = \left(\frac{2}{\pi}\right)^2 \mathcal{N}_l^{\rm den,\gamma} \int dk \, k^2 \, P_{\rm m}(k) \, \mathcal{W}_l^{\rm den}(k_1, k) \, \mathcal{W}_l^\gamma(k_2, k)
\end{equation}

The cross-spectrum kernel combines one density and one lensing transform:
\begin{equation}
\label{eq:window_den_gam_fb}
\mathcal{W}_l^{\rm den}(k_1, k) \mathcal{W}_l^\gamma(k_2, k) = \mathcal{W}_l^{\rm den}(k_1, k) \int_0^\infty d\chi' \, W_l^\gamma(k_2, \chi') \, j_l(k_1 \chi')
\end{equation}

\textbf{Normalization}: The cross-spectrum is often not separately normalized:
\begin{equation}
\mathcal{N}_l^{\rm den,\gamma} = 1
\end{equation}

\textbf{Physical interpretation}: This cross-spectrum couples the local galaxy distribution (at $k_1$) to the integrated matter perturbations (at $k_2$). It is sensitive to growth rate and is particularly valuable for cosmological constraints.

\subsection{Detailed Substitution Procedure: Density × Density Example}

To illustrate the full procedure, we explicitly work out the density × density case.

\subsubsection{Step 1: Define the Observable Kernels}

From Eq.~\ref{eq:kernel_den}, the density kernel is:
\begin{equation}
W^{\rm den}_l(k, z) = \frac{H(z)}{c} \, b(z) \, n(z) \, D(z) \, j_l[k r(z)]
\end{equation}

We abbreviate this as $W(k, z)$ for clarity in this subsection.

\subsubsection{Step 2: Construct the Radial Window}

The window function weighted by a spherical Bessel function is:
\begin{equation}
\widetilde{W}_\chi(k, \chi) = W(k, z(\chi)) \, j_l(k \chi) = \frac{H(z(\chi))}{c} \, b(z(\chi)) \, n(z(\chi)) \, D(z(\chi)) \, j_l[k r(z(\chi))] \, j_l(k \chi)
\end{equation}

Note: Since $\chi = r(z)$, we have $r(z(\chi)) = \chi$, simplifying the first Bessel argument:
\begin{equation}
\widetilde{W}_\chi(k, \chi) = \frac{H(\chi)}{c} \, b(\chi) \, n(\chi) \, D(\chi) \, j_l(k \chi) \, j_l(k \chi) = \frac{H(\chi)}{c} \, b(\chi) \, n(\chi) \, D(\chi) \, [j_l(k \chi)]^2
\end{equation}

\subsubsection{Step 3: Fourier-Bessel Transform (First Observer Wavenumber $k_1$)}

The first Fourier-Bessel transform projects the window function onto the $k_1$ wavenumber:
\begin{equation}
\mathcal{W}_l^{(1)}(k_1, k) = \int_0^\infty d\chi \, \widetilde{W}_\chi(k, \chi) \, j_l(k_1 \chi)
\end{equation}

Substituting the explicit form:
\begin{equation}
\label{eq:step3_den_den}
\mathcal{W}_l^{(1)}(k_1, k) = \int_0^\infty d\chi \, \frac{H(\chi)}{c} \, b(\chi) \, n(\chi) \, D(\chi) \, [j_l(k \chi)]^2 \, j_l(k_1 \chi)
\end{equation}

This integral couples the observer wavenumber $k_1$ to the matter wavenumber $k$ via a triple-product of Bessel functions.

\subsubsection{Step 4: Second Observer Wavenumber ($k_2$)}

The power spectrum couples two independent Fourier-Bessel transforms for wavenumbers $k_1$ and $k_2$:
\begin{equation}
\mathcal{W}_l^{\rm den}(k_1, k) \, \mathcal{W}_l^{\rm den}(k_2, k) = \left[\int_0^\infty d\chi \, \widetilde{W}_\chi(k, \chi) \, j_l(k_1 \chi)\right] \left[\int_0^\infty d\chi' \, \widetilde{W}_\chi(k, \chi') \, j_l(k_2 \chi')\right]
\end{equation}

\subsubsection{Step 5: Integrate Over Matter Wavenumber}

The final power spectrum integrates over all matter wavenumbers, weighted by $k^2 P_{\rm m}(k)$:
\begin{equation}
\label{eq:final_den_den_explicit}
S_l^{\rm den,den}(k_1, k_2) = \left(\frac{2}{\pi}\right)^2 \int_0^\infty dk \, k^2 \, P_{\rm m}(k) \left[\int_0^\infty d\chi \, \frac{H(\chi)}{c} \, b(\chi) \, n(\chi) \, D(\chi) \, [j_l(k \chi)]^2 \, j_l(k_1 \chi)\right]
\end{equation}
\begin{equation}
\quad \times \left[\int_0^\infty d\chi' \, \frac{H(\chi')}{c} \, b(\chi') \, n(\chi') \, D(\chi') \, [j_l(k \chi')]^2 \, j_l(k_2 \chi')\right]
\end{equation}

Rearranging the order of integration (exchanging $d\chi \, d\chi' \, dk$ for $dk \, d\chi \, d\chi'$):
\begin{equation}
\label{eq:final_den_den_reordered}
S_l^{\rm den,den}(k_1, k_2) = \left(\frac{2}{\pi}\right)^2 \int_0^\infty d\chi \, \frac{H(\chi)}{c} \, b(\chi) \, n(\chi) \, D(\chi) \, j_l(k_1 \chi) \int_0^\infty d\chi' \, \frac{H(\chi')}{c} \, b(\chi') \, n(\chi') \, D(\chi') \, j_l(k_2 \chi')
\end{equation}
\begin{equation}
\quad \times \int_0^\infty dk \, k^2 \, P_{\rm m}(k) \, [j_l(k \chi)]^2 \, [j_l(k \chi')]^2
\end{equation}

The innermost integral over $k$ represents the 1D power spectrum projection onto the Bessel basis pairs at distances $\chi$ and $\chi'$. This can be evaluated numerically or approximated depending on the accuracy requirements.

\subsection{Comparison with Limber Approximation}

In the Limber approximation, the steep $k$-dependence of $P_{\rm m}(k)$ is used to approximate the Bessel integrals. The approximation assumes that the main contribution comes from $k \sim l/\chi$, leading to:
\begin{equation}
\label{eq:limber_approx}
\int_0^\infty dk \, k^2 \, P_{\rm m}(k) \, j_l(k \chi) j_l(k \chi') \approx \frac{\pi}{2} \, \frac{1}{\chi^2} \, P_{\rm m}(l/\chi) \, \delta_{\chi \chi'}
\end{equation}

This approximation:
\begin{itemize}
  \item Enforces $\chi = \chi'$ (large-$\ell$ limit: only diagonal elements contribute)
  \item Replaces the Bessel integral with a simple power spectrum evaluation at $k = l/\chi$
  \item Reduces 3D integrals to 1D redshift integrals
  \item Fails at low $\ell$ and small-scale separations
\end{itemize}

The full SFB treatment retains all off-diagonal terms and is valid across all multipole orders.

\subsection{Implementation Strategy: Numerical Quadrature}

\subsubsection{Gauss-Legendre Quadrature for Redshift Integrals}

The Fourier-Bessel transforms are typically computed via Gauss-Legendre quadrature:
\begin{equation}
\mathcal{W}_l^A(k_1, k) \approx \sum_{i=1}^{N_\chi} w_i \, W_l^A(k, z_i) \, j_l(k \chi_i) \, j_l(k_1 \chi_i)
\end{equation}

where $(w_i, \chi_i)$ are the quadrature weights and nodes. The number of nodes $N_\chi$ is chosen to achieve target accuracy.

\subsubsection{Adaptive Quadrature for Wavenumber Integral}

The outer $k$-integral is typically handled with adaptive schemes:
\begin{equation}
S_l^{AB}(k_1, k_2) \approx \sum_j w_j^{(k)} \, (k_j)^2 \, P^{AB}(k_j) \, \mathcal{W}_l^A(k_1, k_j) \, \mathcal{W}_l^B(k_2, k_j)
\end{equation}

Commonly used methods include QuadGK (Gauss-Kronrod with extrapolation) for smooth power spectra.

\subsubsection{Caching for Performance}

For efficiency, the Fourier-Bessel transforms $\mathcal{W}_l^A(k_1, k)$ and $\mathcal{W}_l^B(k_2, k)$ are pre-computed on grids and cached, since they depend only on the observable kernels and not on the power spectrum details.

\subsection{Summary: Complete Workflow}

The complete procedure for inserting observable kernels into the 3D power spectrum is:

\begin{enumerate}
  \item \textbf{Choose observables $A$ and $B$}: Select from density, lensing, RSD, or magnification.
  \item \textbf{Define kernels}: Extract $W_l^A(k, z)$ and $W_l^B(k, z)$ from Eq.~\ref{obs}.
  \item \textbf{Construct redshift windows}: Assemble $\widetilde{W}^A(k, \chi) = W_l^A(k, z(\chi)) \, j_l(k \chi)$.
  \item \textbf{Fourier-Bessel transform}: Compute $\mathcal{W}_l^A(k_1, k)$ and $\mathcal{W}_l^B(k_2, k)$ using quadrature.
  \item \textbf{Integrate over matter wavenumber}: Compute the $k$-integral with adaptive quadrature.
  \item \textbf{Output}: The result is the 3D power spectrum $S_l^{AB}(k_1, k_2)$ with full dependence on both observer wavenumbers.
\end{enumerate}

This procedure is exact (within numerical accuracy) and captures all spherical harmonic and redshift-distance effects beyond the Limber approximation.
